\subsection{Ampliaciones técnicas de los capítulos V y VII}
\label{anx:ampliaciones}

En este anexo se reúne el detalle técnico que los capítulos V y VII resumen. Cada apartado indica la sección que lo remite.

\subsubsection{Asignación de lenguajes: el perfilado y la cota de Amdahl}

La hipótesis de partida sostenía que los componentes intensivos en cómputo debían
reescribirse en un lenguaje compilado. Antes de comprometer esfuerzo se perfiló una lista
escolar de veinte ítems, la mediana del corpus real, contra un catálogo de 3.187 productos.
Las cifras de esta sección son reportadas: las registra la documentación del código fuente,
que remite a un informe de perfilado que no forma parte de la entrega, y el código del grafo
perfilado ya no existe en el sistema construido. Sus salidas originales no se examinaron,
aunque la aritmética entre las cifras es coherente.

El resultado refutó la hipótesis por una vía inesperada. El 96,3\% de la latencia
corresponde, en efecto, a cómputo en el intérprete de Python, pero el 95,3\% del total lo
produce una sola línea, que reconstruye un grafo de sustituibilidad completo en cada ítem
procesado. Construir el grafo cuesta 1.673~ms y consultarlo, 0,137~ms: cuatro órdenes de
magnitud de diferencia. La Tabla~\ref{tab:latencia} presenta el desglose.

% Tabla tab:latencia movida al Anexo A.3 (tablas del cuerpo)


Según la ley de Amdahl \citep{amdahl_validity_1967}, sobre el código tal como estaba y con
una fracción paralelizable $P = 0{,}9633$, una aceleración de $50\times$ produciría una
reducción del 94,4\%, cifra que en apariencia justificaría la reescritura. No la justifica,
porque el mismo efecto se obtiene almacenando el grafo en caché. Corregido ese defecto, la
latencia total desciende a 1.657~ms y la fracción de cómputo en Python cae a
$P = 0{,}2219$, con lo que la reducción alcanzable por cambio de lenguaje queda entre el
20,0\% y el 21,7\%. Las fuentes no declaran las aceleraciones supuestas; esos dos valores
corresponden a $10\times$ y $50\times$ en $P\,(1 - 1/s)$. La cota se refiere al pipeline de
resolución, no al servicio de comercio, cuyo desempeño no tiene medición (ADR 0003).

Esa cota superior es además generosa. El 98,5\% de la fracción residual corresponde al
optimizador NSGA-II \citep{deb_fast_2002} implementado sobre pymoo \citep{blank_pymoo_2020},
biblioteca de la que no se documentó un equivalente maduro en Go. Migrar exigiría
reimplementar el algoritmo, el operador de selección por pseudopesos y el cálculo de
hipervolumen \citep{zitzler_multiobjective_1999}, es decir, reconstruir una biblioteca de
investigación validada para recuperar 362~ms, el 1,03\% del total.

Los seis desplegables del sistema construido están, por tanto, en Python. El único
componente con un argumento técnico para un lenguaje compilado era la ingesta, por su perfil
de concurrencia de red, y el catálogo curado la reemplazó (Sección~\ref{sec:consolidacion}).
La estructura sigue siendo políglota: el manifiesto declara el lenguaje de cada servicio y
una prueba lo verifica, de modo que incorporar un componente en otro lenguaje es una adición
y no una reestructuración.

\subsubsection{Las doce pruebas de fronteras}

Un manifiesto declara, para cada servicio, su lenguaje, su base de datos,
las tablas de las que es dueño y las referencias que cruzan hacia otros dominios, y doce
pruebas lo verifican en cada cambio. Cinco revisan el código: que ningún servicio importe
código de otro, declare una tabla de un dominio ajeno ni mencione la base de datos de otro;
que toda referencia a otro dominio esté declarada, y que el lenguaje coincida con el
declarado. Las otras cuatro revisan el propio manifiesto: que cuadre con el esquema
objetivo, que sus referencias cruzadas sean coherentes, que ninguna tabla tenga dos dueños y
que el grafo de dependencias entre servicios sea acíclico. Dos más, añadidas con las
colas, comprueban que toda lectura de la base de otro servicio esté declarada, exista y sirva
solo de respaldo sin RabbitMQ, y la última, añadida con Redis, que solo la biblioteca común de
los servicios hable con Redis. Los registros de esas
ejecuciones no forman parte de la entrega, así que aquí se describe lo que las pruebas
comprueban y no su historial de resultados.

La justificación de este mecanismo proviene del propio proyecto. El defecto más grave que
encontró el rediseño fue que la tabla de productos del catálogo estaba declarada como
maestra en dos servicios a la vez. No se produjo por falta de criterio, sino porque nada
verificaba el criterio de forma continua, y esa falla de la disciplina humana sobre este
mismo código basta para no depender solo de ella.

\subsubsection{Auditoría de seguridad: defectos y estado}


El de mayor alcance estaba en las políticas de seguridad a nivel de fila de las tablas con datos de
estudiantes. La condición base concedía visibilidad de la tabla completa cuando el contexto de
usuario no estaba establecido, de modo que el error más probable en un sistema distribuido, que un
servicio olvide fijar ese contexto, otorgaba acceso irrestricto. Una segunda cláusula daba acceso
completo al rol de administrador comparando solo el identificador del usuario, sin ancla de ámbito,
y ese defecto motivó el modelo de autorización con ámbito del capítulo de diseño. El ADR 0015 lo
registra en 16 tablas de cinco servicios y advierte que corregirlo va antes que su propia decisión;
las migraciones posteriores conservaban la apertura, sin pruebas de regresión. Entre los demás, el
manejador genérico de errores de validación devolvía al cliente el valor rechazado, contraseña
incluida, y la supresión de datos se marcaba como completada sin confirmación de los servicios.


El sistema construido no tiene políticas de seguridad a nivel de fila. El ámbito se controla en la
aplicación, con la cadena de contención de identidad, que falla cerrada
(Sección~\ref{sec:autorizacion}); las cifras de 152 y 22 son las de la auditoría, no una medición
propia. De los catorce defectos del Anexo~M, dos están corregidos: el valor rechazado ya no vuelve
al cliente, y el correo del acudiente ya no queda en un registro que la supresión no alcanza. El
texto de la lista sin depurar se corrigió para la lista oficial, pero sigue en el documento que
carga el acudiente; la depuración del nombre se sustituyó por la minimización
(Sección~\ref{sec:percepcion-dev}), salvo en ese mismo documento. Los tres defectos de la supresión
cambiaron de forma, pero la supresión sigue sin mecanismo: el estado de propagación existe en el
modelo, sin publicador ni consumidor del evento.

Cuatro siguen vigentes: la gestión de cuentas no deja registro de auditoría, el consentimiento por
finalidad no se escribe ni se consulta, la respuesta de la importación de listas guarda el extracto
completo del documento y ninguna tabla de auditoría tiene una política de purga. Otro no tiene
efecto mientras no se escriba el registro de auditoría. Uno es parcial: las preferencias y los pesos
del acudiente están cifrados, pero la razón de la recomendación queda en claro y no se borra. El de
la anonimización no es evaluable, porque no hay ese proceso.

Un defecto de comercio quedó corregido en la versión evaluada. La consulta de los pedidos del
acudiente toma el usuario de la sesión, validada con identidad, e ignora el que llegue en la
consulta; el detalle, el pago y la cancelación de un pedido propio también exigen sesión, y un
pedido ajeno responde 404. La puerta de enlace sigue sin exigir sesión.

En el sistema construido siguen abiertos otros defectos. La consulta de contactos del servicio de identidad devuelve el correo de
cualquier usuario activo sin autenticación, y las rutas de los paneles de catálogo, inteligencia y
operación no exigen sesión ni rol. Cuatro caminos crean matrículas, y solo uno, el registro de un
hijo desde la cuenta del acudiente, escribe el consentimiento del menor. El celular del acudiente se
guarda sin cifrar en la base de identidad.

Las pruebas propias del sistema no detectaron ninguno de estos defectos. Medir cuánto código
ejecutan no dice si buscan las condiciones adversas correctas, y varias comprueban garantías leyendo
el código fuente en lugar de ejecutar el comportamiento. La verificación de la versión evaluada
(Sección~\ref{sec:calidad}) los confirmó y encontró otros. Un kit se puede comprar sin sesión, aun
mientras espera revisión, y varias rutas internas aceptan llamadas sin verificar quién llama. Las
cinco bases de datos comparten un mismo rol de acceso, de modo que las credenciales de cualquier
servicio leen los usuarios de identidad. La página del documento cargado sale sin depurar hacia el
proveedor del modelo (\mbox{PE-06}), y el nombre ficticio de las pruebas quedó guardado. El retracto existe
en comercio, pero ninguna aplicación lo ofrece y la puerta de enlace no lo expone; además, se acepta
pasado el plazo legal y no libera la reserva. Un pedido de 100.000 unidades fue aceptado para pago.
La respuesta de la puerta de enlace expone direcciones internas, el mensaje de un trabajo fallido
contiene la consulta SQL con sus parámetros y, sin la llave de cifrado, el cifrador guarda en claro
sin avisar.

La exposición de estos defectos tiene una consecuencia de alcance. Las pruebas con usuarios usaron
estudiantes ficticios y una pasarela de pruebas, de modo que ningún dato real de menores ni ningún
cobro real quedó expuesto. Cerrar los defectos que siguen abiertos antes de recibir familias reales
forma parte del endurecimiento productivo que se delimita en la
Sección~\ref{sec:trabajo-futuro}. El retracto, que la Ley 1480 de 2011 reconoce al comprador, se
completó en la aplicación del acudiente y en la puerta de enlace\pendienteEvidencia{retracto ejercible desde la aplicacion, con liberacion de la reserva y plazo legal}.

\subsubsection{Estado de recuperación por fase de CRISP-DM}

La Tabla~\ref{tab:rag-fases} resume, por fase de CRISP-DM, las 84 estrategias de recuperación evaluadas.

\begin{table}[htbp]
\centering
\caption{Estrategias de recuperación por fase de CRISP-DM y estado en el sistema construido}
\label{tab:rag-fases}
\footnotesize
\begin{tabular}{@{}lrrrrp{6.2cm}@{}}
\toprule
\textbf{Fase} & \textbf{Impl.} & \textbf{Prot.} & \textbf{Desc.} & \textbf{Otras} & \textbf{Lo que usa el sistema construido} \\
\midrule
Comprensión del negocio & 1 & 0 & 1 & 5 & Sugerencias ancladas al catálogo, sin generación de texto \\
Comprensión de los datos & 2 & 1 & 0 & 0 & Catálogo curado desde la demanda; 174 listas de 17 colegios \\
Preparación de los datos & 13 & 6 & 1 & 4 & Metadatos y niveles del catálogo; percepción con Gemini validada por esquema; limpieza, corrección, sinónimos y atributos del renglón \\
Modelado & 11 & 9 & 5 & 10 & gemini-embedding-001; fusión aditiva del puntaje léxico y el coseno; filtros de metadatos; hasta tres opciones, una por nivel \\
Evaluación & 1 & 3 & 0 & 0 & $\tau=0{,}75$ con abstención y revisión de la lista completa \\
Despliegue & 6 & 1 & 1 & 3 & pgvector en la base del catálogo; caché de vectores en Redis; índice de conceptos en memoria; trabajos durables con espera humana \\
\midrule
Total & 34 & 20 & 8 & 22 & \\
\bottomrule
\end{tabular}
\par\smallskip
{\footnotesize Impl.: implementada; Prot.: solo en prototipos; Desc.: evaluada y descartada; Otras: considerada sin evaluar, no considerada o no aplica al dominio.}
\end{table}


\subsection*{Ampliación de la Sección VII: verificación de la versión del 8 de octubre}

Las fases 0 a 6 corrieron sobre la versión de integración del 8 de octubre de 2026, cuyo código de
producto es el mismo de la versión evaluada salvo la integración continua y dos valores por omisión
de la base de datos. El primer nivel de las fases 7 a 9 corrió sobre la versión evaluada, en una
sola corrida, en una máquina virtual de 4 procesadores y 8~GiB y con la percepción simulada; el
simulador atendió 2.342 percepciones, sin ninguna llamada real al modelo. Con las pruebas que añadió
la verificación pasaban 3.911, otras 129 eran fallas esperadas que documentan defectos y la
cobertura subía a 94,32\,\%. La integración continua ejecuta las pruebas pero no mide la cobertura,
y su estado en esa versión no se comprobó. De las 60 filas de la matriz que liga cada meta con su
ensayo se ejecutaron 36; de las otras 24, 13 corresponden a los niveles 2 y 3, 2 dependen de un
hallazgo ya anticipado, 8 requieren componentes que el sistema construido incorpora después y 1 es
el protocolo con usuarios. Los 78 hallazgos abiertos se reparten en 14 de severidad alta, 56 media
y 8 baja. En la telemetría, las trazas no llevan el identificador de correlación, así que una carga
no se puede seguir de extremo a extremo.

\subsection*{Ampliación de la Sección VII: carga, degradación y defectos de seguridad}

El primer nivel de las fases 7 a 9 dejó 18 ensayos, de los que 14 cumplen y 4 no. En la carga, una
rampa hasta 100 usuarios generó 27.894 peticiones sin errores, y el percentil 95 de la consulta del
catálogo fue de 22~ms con un usuario y de 160~ms con cien, por debajo de los 300~ms de la meta. No
hubo ningún rechazo por límite de tasa porque la versión evaluada no tenía limitador, de modo que en
ella esa meta no se cumple; el sistema construido añade uno con Redis, que en la corrida del 2 de noviembre respondió 429 a
30.257 peticiones que superaban el límite, con cinco usuarios a través de la puerta de enlace. En la degradación, siete de diez caídas de un componente cumplieron: los demás servicios
siguieron respondiendo sin errores y se recuperaron en 2,5~s o menos. Con el catálogo colgado, el
disyuntor abrió a los 26,3~s, dentro de los 30~s del escenario ESC-R2, porque cada intento falló
en unos 5~s. Si cada intento agotara su espera, el modelo da 300,8~s desde inteligencia y 150,8~s
desde comercio. El peor caso se acotó con ese modelo y no se midió; se corrige en el endurecimiento productivo
(Sección~\ref{sec:trabajo-futuro}). Pasados los
30~s, con diez llamadas simultáneas, el disyuntor deja pasar las diez en lugar de una sola de
prueba.

Las tres caídas que no cumplieron revelaron tres limitaciones de la versión evaluada: con el
catálogo detenido, subir una lista respondía con un error 500 y no se guardaba; tras caer la base o
el intermediario, los procesos de fondo morían sin reiniciarse; y un mismo pago producía dos avisos.
El sistema construido corrige las tres (ADR 0031, 0032 y 0033). La corrección se verificó el 20 de
octubre de 2026 en una versión de trabajo: la subida responde 202 en 33~ms y se recupera en 2,7~s,
los procesos de fondo se recuperan en 7,1~s y 6,9~s, y el pago produce un solo aviso. Esa
verificación no volvió a correr los demás ensayos, que se repiten en la corrida completa sobre el
sistema construido\pendienteEvidencia{corrida completa de la verificacion sobre el sistema construido}.


La auditoría de los microservicios anteriores encontró defectos de gravedad desigual. El de mayor
alcance estaba en las políticas de seguridad a nivel de fila: la condición base concedía la tabla
completa cuando el contexto de usuario no estaba fijado, de modo que el error más probable en un
sistema distribuido, olvidar fijarlo, otorgaba acceso irrestricto (ADR 0015, 16 tablas de cinco
servicios). El sistema construido no tiene políticas de fila: controla el ámbito en la aplicación,
con la cadena de contención de identidad, que falla cerrada (Sección~\ref{sec:autorizacion}).
De los catorce defectos del Anexo~M, dos están corregidos (el valor rechazado ya no vuelve al
cliente y el correo del acudiente ya no queda en un registro que la supresión no alcanza), uno es
parcial, cuatro siguen vigentes y uno no es evaluable. En comercio, la consulta de pedidos toma el
usuario de la sesión y un pedido ajeno responde 404. Siguen abiertos defectos de acceso (rutas de
catálogo, inteligencia y operación sin sesión ni rol, y un rol de acceso compartido por las cinco
bases) y de privacidad (el nombre del menor dentro del documento cargado y la supresión de datos
sin mecanismo). Las pruebas propias del sistema no detectaron ninguno, porque medir cuánto código
ejecutan no dice si buscan las condiciones adversas correctas; la verificación de la versión
evaluada (Sección~\ref{sec:calidad}) los confirmó. El detalle de cada defecto está en el
Anexo~\ref{anx:ampliaciones}.

